%% ma: L12-B18-0003
%% lop: 12
%% chuong: 6
%% bai: 18
%% dang: 12.18.3
%% loai: TLN
%% muc: VD
%% dap_an: "0,3"
%% nguon: "(NBV)-ĐỀ SỐ 7-MỖI NGÀY 1 ĐỀ THI 2025.pdf (D07, MỖI NGÀY 1 ĐỀ - Đề số 7), Phần III Câu 4"
%% kien_thuc: "Bài 18 (xác suất có điều kiện), công thức cộng xác suất (Lớp 11 Bài 29)"
%% trang_thai: chinh-thuc
%% ghi_chu: "Bỏ ảnh biểu tượng thuốc lá và răng (chỉ minh họa). Đáp án gốc 0,3 đúng."
\begin{ex}
Một tổ chức y tế tiến hành khảo sát mối liên hệ giữa việc hút thuốc và khả năng mắc các bệnh về răng miệng của một cộng đồng người trong một khu vực. Thống kê cho thấy số người hút thuốc chiếm tỉ lệ $30\%$, số người mắc các bệnh về răng miệng chiếm tỉ lệ $50\%$ và số người hoặc hút thuốc, hoặc mắc các bệnh về răng miệng hoặc cả hai chiếm tỉ lệ $65\%$. Chọn ngẫu nhiên một người từ khu vực X. Tính xác suất người đó hút thuốc với điều kiện người đó không mắc các bệnh về răng miệng.
\shortans{0,3}
\loigiai{Gọi $A$: "người hút thuốc", $B$: "người mắc bệnh răng miệng". Có $P(A)=0{,}3$, $P(B)=0{,}5$, $P(A\cup B)=0{,}65$. Theo công thức cộng, $P(A\cap B)=0{,}3+0{,}5-0{,}65=0{,}15$. Suy ra $P(A\cap\overline B)=P(A)-P(A\cap B)=0{,}15$ và $P(\overline B)=0{,}5$. Vậy $P(A\mid\overline B)=\dfrac{P(A\cap\overline B)}{P(\overline B)}=\dfrac{0{,}15}{0{,}5}=0{,}3$.}
\end{ex}
